\scanpage{161}
% archive sheet number: 160
% L1
\uline{Cas des préschémas} loc. noeth.

% L2-5
Si $F$ est cohérent, les $\underline{\operatorname{Ext}}^i_Z(F,G)$ sont quasi-cohérents. On le voit par \unsure{exemple|récurrence} sur \unsure{la} suite spectrale de terme initial $\underline{H}^p_Z(\underline{\operatorname{Ext}}_{\mathcal O_X}(F,G))$, utilisant le fait qu'une extension de qu. coh. est \unsure{itou} \unsure{...}

% L6-8
On a un hom. canonique (\unsure{$Y$|$Z$} fermé)
\[
\left\{\begin{aligned}
1)\quad &\operatorname{Ext}^i_{\deleted{\unsure{Y}}}(X;F,G)\longleftarrow\varinjlim\operatorname{Ext}^i(X;F/{\unsure{\mathcal J|J}}^{n+1}F,G)\\
2)\quad &\underline{\operatorname{Ext}}^i_{\unsure{Y|\gamma}}(F,G)\longleftarrow\varinjlim\underline{\operatorname{Ext}}^i(F/\mathcal J^{\unsure{n+1}}F,G)
\end{aligned}\right.
\]

% L9-11
\uline{Th.} \inserted{Supposons $F$ \uline{cohérent}.} L'hom. 2) est un isom, \deleted{\unsure{$X$ loc. noeth.}}, et 1) est un isom si $X$ est noethérien.

% L12-14
On déduit \deleted{\unsure{de}} le \unsure{deuxième} isom du premier comme d'habitude, par la suite spectrale
\[
H^p(X,\underline{\operatorname{Ext}}^q_Y(F,G))\Longrightarrow\operatorname{Ext}_Y(X;F,G).
\]
% L15-18
L'énoncé est de nature \uline{\unsure{loc.|locale}}, on peut supposer $X$ affine. Pour $i=0$, c'est évident, il reste à \deleted{\unsure{examiner}} prouver que si $G$ provient d'un module \uline{injectif}, on a $\underline{\operatorname{Ext}}^i_Y(F,G)=0$ pour $i>0$.

% L19
Or on regarde $\operatorname{Ext}^p_A(\unsure{M|N},\underline{H}^q_J(N))$ et on voit que si $N$
